\documentclass[11pt]{article}
\usepackage[T1]{fontenc}
\usepackage[utf8]{inputenc}
\usepackage[margin=1in]{geometry}
\usepackage{amsmath,amssymb,amsthm,booktabs,array,microtype}
\usepackage[numbers,sort&compress]{natbib}
\usepackage[colorlinks=true,linkcolor=blue,citecolor=blue,urlcolor=blue]{hyperref}
\newtheorem{proposition}{Proposition}
\newtheorem{theorem}[proposition]{Theorem}
\newtheorem{lemma}[proposition]{Lemma}
\newtheorem{corollary}[proposition]{Corollary}
\DeclareMathOperator{\KL}{KL}
\DeclareMathOperator{\kl}{kl}
\newcommand{\E}{\mathbb E}
\newcommand{\Prb}{\mathbb P}
\title{Theoretical extensions for verified replay:\\optimization, admission, entropy, and survival}
\author{Companion research note for the ModeBench / ReplayMaxRL manuscript}
\date{September 5, 2026}
\begin{document}
\maketitle

The four follow-up questions admit useful theoretical answers with distinct
premises. This note gives complete derivations, identifies their literature
basis, and states what evidence would be needed to apply them to the runs.
These are applications and specializations of established mathematical tools;
no foundational novelty is claimed. The note contains no new training results.

\begin{center}
\small
\begin{tabular}{@{}p{.19\textwidth}p{.36\textwidth}p{.37\textwidth}@{}}
\toprule
Question & Available theoretical result & Required application evidence \\
\midrule
Optimizer & Simultaneous high-probability retention from a controlled stochastic energy budget & Bounds on update bias, noise, smoothness, and preconditioner dependence \\
Admission & Adaptive hazard bounds plus retention across changes in the bank & Actual insertion opportunities, capacity, persistent coefficients, and switching costs \\
Exact entropy & Same full-support optimum as replay under conditional entropy, but different KL geometry & An exact estimator and explicit output alphabet; no transfer from a historical score \\
Per-mode survival & Sharp probability certificates and valid finite-count intervals & Complete-response likelihoods or per-key sampling counts under a specified decoding law \\
\bottomrule
\end{tabular}
\end{center}

\section{An optimizer-aware conditional guarantee}
\label{sec:optimizer}
% Standalone fragment: requires amsmath, amssymb, amsthm, hyperref,
% and a theorem environment. No labels from either manuscript are required.
\paragraph{Stochastic retention with a controlled energy budget.}
The argument combines a standard variable-metric descent estimate
(\href{https://doi.org/10.1137/15M1053141}{Wang et al., 2017, Section 2})
with the nonnegative-supermartingale maximal inequality
(\href{https://doi.org/10.1214/18-PS321}{Howard et al., 2020, Lemma 1}).
The probability consequence for protected exemplars is derived here; it is
not an existing theorem about Adam or language-model replay.

\begin{theorem}[Optimizer-aware conditional exemplar retention]
Fix finitely many complete verified responses $(x_j,e_j)$ with keys $b_j$,
lengths $L_j>0$, and weights $\alpha_j>0$. Let
\[
 R(\theta)=\sum_j\frac{\alpha_j}{L_j}
             [-\log\pi_\theta(e_j\mid x_j)],\qquad
 F(\theta)=R(\theta)-J(\theta),\qquad J(\theta)\le J_{\max}<\infty.
\]
Assume $F$ is differentiable with globally $L_F$-Lipschitz gradient,
and take deterministic $\theta_0$ with finite $F(\theta_0)$. Relative to a filtration
$(\mathcal F_t)$, consider
\[
 \theta_{t+1}=\theta_t-\eta_t H_t(g_t+b_t+\xi_t),\qquad
 g_t=\nabla F(\theta_t).
\]
Here $\theta_t,\eta_t,H_t,b_t$ are $\mathcal F_t$-measurable,
$0\preceq H_t\preceq MI$, $\eta_t>0$, $\eta_tL_FM\le1$, and
$\mathbb E[\xi_t\mid\mathcal F_t]=0$ with finite conditional second moment.
Suppose deterministic numbers $d_t\ge0$ bound
\[
 \frac{\eta_t}{2}\|b_t\|_{H_t}^2
 +\frac{L_F\eta_t^2}{2}
       \mathbb E[\|H_t\xi_t\|_2^2\mid\mathcal F_t]\le d_t,
 \qquad \|u\|_H^2=u^\top Hu.
\]
For $N\in\mathbb N\cup\{\infty\}$ with
$B_N=\sum_{t<N}d_t<\infty$, put
$D_0=F(\theta_0)+J_{\max}\ge0$. For every $0<\delta<1$,
\[
 \Pr\!\left(
  \forall t\le N,\ \forall j:\quad
  p_{\theta_t}(b_j\mid x_j)
  \ge\pi_{\theta_t}(e_j\mid x_j)
  \ge\exp\!\left[-\frac{L_j(D_0+B_N)}{\alpha_j\delta}\right]
 \right)\ge1-\delta.
\]
When $N=\infty$, each protected probability also has an almost surely
positive, potentially random, infimum over all iterations.
\end{theorem}

\begin{proof}
Smoothness, conditional centering of $\xi_t$, and
$H_t^2\preceq MH_t$ give
\begin{align*}
 \mathbb E[F(\theta_{t+1})\mid\mathcal F_t]
 &\le F(\theta_t)-\eta_tg_t^\top H_t(g_t+b_t)
       +\frac{\eta_t}{2}\|g_t+b_t\|_{H_t}^2
       +\frac{L_F\eta_t^2}{2}
          \mathbb E[\|H_t\xi_t\|_2^2\mid\mathcal F_t]\\
 &\le F(\theta_t)-\frac{\eta_t}{2}\|g_t\|_{H_t}^2+d_t.
\end{align*}
The cross term involving $b_t$ cancels on expanding the square.
Since $F+J_{\max}\ge R\ge0$, the adapted process
\[
 X_t=F(\theta_t)+J_{\max}+\sum_{s=t}^{N-1}d_s
\]
is a nonnegative supermartingale; for $N=\infty$ use the infinite tail.
Its initial value is $D_0+B_N$. The maximal inequality yields
$\Pr(\sup_{t\le N}X_t\ge(D_0+B_N)/\delta)\le\delta$.
If $D_0+B_N=0$, nonnegativity instead gives $X_t=0$ almost surely.
Every nonnegative summand of $R(\theta_t)\le X_t$ obeys the same bound,
and exponentiation proves the simultaneous exemplar floors without a
union factor over $j$. Each complete-exemplar event is contained in its
verified-key event. For $N=\infty$, the same maximal inequality, with a
threshold tending to infinity, gives $\sup_tX_t<\infty$ almost surely.
\end{proof}

\paragraph{Finite and infinite horizons are different.}
Unbiased bounded-variance updates have $d_t=O(\eta_t^2)$; square-summable
steps therefore suffice for the infinite-horizon energy budget.
A persistent bias instead costs $O(\eta_t\|b_t\|^2)$.
Constant steps and persistent variance give only the stated finite-horizon
bound unless additional structure controls accumulated noise.
No positive lower eigenvalue of $H_t$ is needed for retention alone;
this theorem does not assert parameter convergence or optimality.
The likelihood must include the complete response and termination event
under the sampling law whose key probability is bounded.

\paragraph{Scope for adaptive implementations.}
A current-gradient Adam second moment is not predictable with respect to
that gradient's noise, and momentum, clipping, bank changes and cyclic
prompt scheduling need separate error or energy control. Positive replay
weight and a small nominal learning rate do not verify these assumptions.
\href{https://arxiv.org/html/1904.09237v1}{Reddi et al. (2018)} provide
counterexamples to unrestricted Adam convergence; this does not itself
prove replay failure in the present application. The bound above is
conditional and can be numerically weak, particularly for long exemplars
or small replay weights.


\section{From actual admission to persistent protection}
\label{sec:discovery}
Conditional Borel--Cantelli supplies the standard probability background
\cite[Theorem~4.3.4]{durrett2019}. Controlling energy across switches has a
long history in hybrid-system analysis \cite{branicky1998}; the elementary
positive-jump accounting below is proved directly.
Archive methods such as Go-Explore separate discovery, return, and policy
learning \cite{ecoffet2021}. That is a conceptual precedent, not a theorem
about the present response bank.
% Standalone fragment: requires amsmath, amssymb, amsthm, hyperref and proposition/corollary environments.
% Suggested citations are specified with verified metadata in discovery.md.
% This fragment has no dependencies on the main manuscript's theorem labels.

\paragraph{Discovery requires admission opportunities.}
The following is a standard conditional-hazard argument, closely related to
conditional Borel--Cantelli
(\href{https://sites.math.duke.edu/~rtd/PTE/PTE5_011119.pdf}{Durrett, 2019,
Theorem~4.3.4}), applied to insertion of a verified key into memory.
Its hazard includes scheduling, validation, capacity, and candidate selection;
a positive probability of generating a key alone is insufficient.

\begin{proposition}[Conditional admission hazards]
Fix a key $b$ not initially banked. Let $(\mathcal F_n)_{n\ge0}$ describe the
complete training history after each opportunity, and let $\tau_b$ be its first
admission time, with $\tau_b=\infty$ if it is never admitted. Define the
predictable stopped hazard
\[
 h_{b,n}=\Pr(\tau_b=n\mid\mathcal F_{n-1}),\qquad
 H_{b,N}=\sum_{n=1}^N h_{b,n}.
\]
For every $a\ge0$,
\[
 \Pr(\tau_b>N,\ H_{b,N}\ge a)\le e^{-a}.
\]
In particular,
$\Pr(\tau_b=\infty,\ H_{b,\infty}=\infty)=0$.
If deterministic numbers $\underline h_{b,n}\in[0,1]$ satisfy
$h_{b,n}\ge\underline h_{b,n}$ on $\{\tau_b\ge n\}$, then
\[
 \Pr(\tau_b>N)\le\prod_{n=1}^N(1-\underline h_{b,n})
 \le\exp\!\left(-\sum_{n=1}^N\underline h_{b,n}\right).
\]
For a finite target set $S$, the probability that some target key has not been
admitted by $N$ is at most the sum of these bounds over $b\in S$.
\end{proposition}
\begin{proof}
Let $Z_N=\mathbf1\{\tau_b>N\}\exp(H_{b,N})$, with $Z_0=1$.
The event of an earlier admission is known at time $n-1$, and $h_{b,n}=0$
on that event. On survival through $n-1$,
\[
 \mathbb E[Z_n\mid\mathcal F_{n-1}]
 =Z_{n-1}e^{h_{b,n}}(1-h_{b,n})\le Z_{n-1}.
\]
Thus $Z_n$ is a nonnegative supermartingale. Markov's inequality gives the
finite-hazard bound. Fatou's lemma excludes a positive-probability event on
which survival continues forever while $H_{b,N}\to\infty$, since then
$Z_N\to\infty$. Under deterministic hazard floors, conditioning gives
$\Pr(\tau_b>n)\le(1-\underline h_{b,n})\Pr(\tau_b>n-1)$;
iteration and the union bound finish the proof.
\end{proof}

Divergent cumulative hazard is a substantive premise. Independent
opportunities with admission probabilities $h_n=(n+1)^{-2}>0$ satisfy
\[
 \Pr(\tau=\infty)=\prod_{n=1}^\infty
 \left(1-\frac1{(n+1)^2}\right)=\frac12.
\]
Likewise, once a non-evicting capacity-$K$ bank is full, every unbanked key has
zero admission hazard. These bounds establish coverage only for keys that
continue to have admissible opportunities; they do not establish such
opportunities for the implemented policy or for supports larger than capacity.

\paragraph{Retention across bank changes.}
An admitted key also needs persistent loss weight. The next elementary energy
accounting result allows bank changes while explicitly charging their effect
on the objective. It is not an assumption that different bank objectives
share a single unchanged potential. This follows the standard switched-energy
viewpoint (\href{https://people.ece.ubc.ca/moishi/eece571m/articles/Branicky98.pdf}{Branicky, 1998});
the specific inequality below is proved directly.

\begin{proposition}[Switching-energy retention]
Let $t_0<t_1<\cdots$ be switching times with finitely many switches on every
bounded time interval. On interval $r$, use a finite collection of verified
complete-response exemplars, with nonnegative coefficients $d_{rj}$ and
\[
 F_r(\theta)=\sum_{j:d_{rj}>0}d_{rj}\ell_j(\theta)-J(\theta),
 \qquad \ell_j(\theta)=-\log\pi_\theta(e_j\mid x_j),
 \qquad J(\theta)\le J_{\max}<\infty.
\]
Zero coefficients denote absent exemplars and are omitted from the sum.
Suppose initial energy is finite, parameters are continuous
at switches, $F_r$ is nonincreasing within its interval, and the positive
energy jumps obey
\[
 A:=\sum_{r\ge1}
 \left[F_r(\theta(t_r))-F_{r-1}(\theta(t_r))\right]_+<\infty.
\]
If exemplar $j$ is protected from interval $r_j$ onward and
$d_{rj}\ge\underline d_j>0$ for every $r\ge r_j$, then, with
$D=F_0(\theta(t_0))+A+J_{\max}<\infty$,
\[
 p_\theta(b_j\mid x_j)\ge\pi_\theta(e_j\mid x_j)
 \ge\exp(-D/\underline d_j)
 \qquad(t\ge t_{r_j}).
\]
The likelihood must represent the complete verified response under the law
whose key probability is bounded. Mean-token losses have
$d_{rj}=\alpha_{rj}/L_j$ for fixed exemplar length $L_j>0$.
\end{proposition}
\begin{proof}
Induction over the switches gives
$F_r(\theta(t))\le F_0(\theta(t_0))+A$. Therefore
$\sum_jd_{rj}\ell_j(\theta(t))\le D$.
Each summand is nonnegative, so a protected exemplar has
$\ell_j(\theta(t))\le D/\underline d_j$.
Its complete response event is contained in its verified-key event.
\end{proof}

If old coefficients are preserved when new exemplars are added, a switch
costs exactly the summed weighted surprisal of the new exemplars at admission.
For uniform categorical replay with fixed dose $\rho$, adding one key to a
size-$k$ bank instead changes the potential by
\[
 \Delta F=\frac{\rho}{k+1}
 \bigl[-\log p_{\rm new}-R_{\rm old}\bigr].
\]
Old coefficients decrease from $\rho/k$ to $\rho/(k+1)$, but remain at least
$\rho/K$ when capacity is $K$. Finitely many such admissions, with positive
admission-time probabilities and no eviction, therefore have finite total
switching cost in the exact descent model. Recurrent scheduling, bounded
priority multipliers, or a finite-capacity buffer alone do not establish
within-interval descent or a summable jump budget for arbitrary optimization.


\begin{corollary}[Conditional discovery followed by retention]
Suppose a prompt has a finite correct support $\mathcal C$ of size $m$ no
larger than bank capacity. Admissions are non-evicting, and every missing key
$b$ has conditional admission hazard at least $\epsilon r_b$ at every
opportunity, for constants $\epsilon>0$ and $r_b>0$ with
$\epsilon r_b\le1$. Assume also the switching-energy conditions above,
including a finite positive-jump budget and persistent positive coefficients
for the admitted complete exemplars. Then the bank reaches full correct
coverage in finitely many opportunities almost surely, and thereafter every correct key has
a uniform-in-time positive probability floor. At horizon $N$,
\[
 \Pr(\text{full coverage by }N)
 \ge 1-\sum_{b\in\mathcal C}e^{-N\epsilon r_b}.
\]
If $r_b\ge\gamma>0$ for every correct key, this is at least
$1-m e^{-N\epsilon\gamma}$.
\end{corollary}
\begin{proof}
Apply the admission-hazard bound to each initially missing key and take a
union bound. The failure bound tends to zero; finite support and non-eviction
therefore imply a finite almost-sure opportunity index of simultaneous full coverage.
The switching-energy bound protects each admitted complete exemplar and hence
its key, and applies to all keys after that time.
\end{proof}

A fixed full-support proposal mixture supplies such a hazard floor if it is
queried at every counted opportunity and insertion is guaranteed whenever the
key appears. This is a sufficient design condition, not a property established
for the current temperature schedule, bounded search attempts, or adaptive
priority controller. The corollary guarantees retention, not uniform key
probabilities or convergence of the neural optimizer.


\section{Exact entropy and replay: the same target can have different protection}
\label{sec:entropy-geometry}
The KL-direction distinction is established in the label-smoothing literature
\cite[Section~3.2]{pereyra2017}. The corresponding log-barrier regularizer is
also analyzed for policy gradient \cite[Section~5.2]{agarwal2021}. Our
specialization concerns a target restricted to verified modes, rather than
uniform regularization over every action. A stored bank supplies a target
that can be scored even when fresh sampling misses it.

Let $\mathcal C$ contain $m\ge2$ correct categories, let
$P=\sum_{c\in\mathcal C}p_c>0$, $q_c=p_c/P$, and let $u_c=1/m$.
Write $H(q)=-\sum_cq_c\log q_c$, with $0\log0=0$.
\begin{proposition}[Full-support objective comparison]
\label{prop:entropy-objectives}
For a full correct-mode bank with a uniform categorical target,
\[
 H(q)=\log m-\KL(q\|u),\qquad
 R_u(p):=-\frac1m\sum_{c\in\mathcal C}\log p_c
 =\log m-\log P+\KL(u\|q).
\]
If $\Psi$ is continuous and strictly increasing on $[0,1]$, then for any
$\beta,\rho>0$, both
\[
 \Psi(P)+\beta H(q)
 \quad\hbox{and}\quad
 \Psi(P)-\rho R_u(p)
\]
have the unique maximizing distribution $P=1,q=u$ on the closed simplex
(with the conditional-entropy objective considered on $P>0$).
This is a statement about global maximizers, not about convergence of
stochastic updates or conditional-entropy gradient flow.
\end{proposition}
\begin{proof}
Substitute $p_c=Pq_c$ into $R_u$, and expand both KL divergences.
Since $H(q)\le\log m$, equality only at $q=u$, maximizing the first
objective selects $q=u$ and then $P=1$. The second is
$\Psi(P)+\rho\log P-\rho\log m-\rho\KL(u\|q)$; its first two terms
increase strictly with $P$, while the last is maximized only at $q=u$.
At that point there is zero incorrect mass, so the distribution is unique.
\end{proof}

Full-output entropy is a different comparator. For a linear correctness
potential $aP$, exact maximization of $aP+\beta H(p)$ gives
\[
 p_i^*=\frac{\exp(a\mathbf1\{i\in\mathcal C\}/\beta)}
 {m\exp(a/\beta)+n},
\]
where $n\ge1$ is the number of incorrect categories. This is the standard
entropy/conjugacy calculation \cite[Section~2.2]{geist2019}. It has
$P^*<1$, whereas exact conditional entropy has no such loss of correctness
at its global optimum. The experimental historical semantic score is not
an exact gradient of either objective.

Existing discrete-time theory is stronger than merely finding an optimum:
Mei et al. prove positive action-probability floors and geometric convergence
for their exact tabular entropy-regularized policy-gradient update
\cite[Lemma~13 and Theorem~5]{mei2020}. This provides relevant prior theory
for an exact finite categorical comparator, not for the historical predictor,
clipped PPO, or AdamW. For the paper's Dr.GRPO reduction the reward vector is
$a\mathbf1_{\mathcal C}$ with $a=(G-1)/G\in(0,1)$, so their bandit reward
assumption applies. Their linear-reward theorem does not directly cover the
nonlinear finite-group MaxRL potential.

The converse directions of KL explain a certificate difference:
$\KL(u\|q)$ diverges if any protected coordinate vanishes, while
$\KL(q\|u)$ remains finite. For example, losing one mode and distributing
mass uniformly over the rest leaves $H(q)=\log(m-1)$. Section~\ref{sec:survival}
quantifies this sharply. This level-set fact does not imply that exact
entropy dynamics necessarily lose a mode.


\section{Sharper per-mode certificates and measurement limits}
\label{sec:survival}
% Generated from survival.tex for the combined note.


This section applies standard relative-entropy data processing, binary-divergence
inversion, and binomial inference to replay retention. The information-theory
facts are established results \cite{vanerven2014}; binary-KL inversion is also a
standard tool in statistical learning \cite{garivier2011}. The propositions
below give self-contained specializations, not claims of new foundational
inequalities. They distinguish deterministic probability certificates from
sampling confidence and from training-trajectory guarantees. All logarithms are
natural. No claim here establishes energy descent for a stochastic neural
optimizer.

\subsection{A sharp certificate from normalized replay cross entropy}
For probability vectors, write $H(w)=-\sum_iw_i\log w_i$ and
\[
 \kl(a,r)=a\log\frac a r+(1-a)\log\frac{1-a}{1-r},
\]
with the usual continuous endpoint conventions. For $0<a<1$ and $D\ge0$,
define $\ell_-(a,D)$ as the unique $r\in(0,a]$ satisfying
$\kl(a,r)=D$. Define $\ell_-(1,D)=e^{-D}$.

\begin{proposition}[Sharp replay certificate]
Let $w_b>0$, $\sum_{b\in B}w_b=1$, and let $p$ be a distribution on an
alphabet containing the bank $B$. Suppose
\[
 R_w(p)=-\sum_{b\in B}w_b\log p_b\le C<\infty,
 \qquad D=C-H(w).
\]
Then $D\ge0$ and every banked coordinate satisfies
\[
 p_b\ge\ell_-(w_b,D)>0.
\]
The bound is sharp on the closed probability simplex. Furthermore,
\[
 p(B)\ge e^{-D},\qquad
 p_b\ge\max\left\{0,w_b-\sqrt{D/2}\right\}.
\]
For $w_b<1$ a simple positive lower bound, weaker than binary inversion, is
$\exp[-(D+h_2(w_b))/w_b]$, where $h_2(a)=H((a,1-a))$.
\end{proposition}

\begin{proof}
Extend $w$ by zero off the bank, writing $\bar w$. Normalization gives
$\KL(\bar w\|p)=R_w(p)-H(w)\le D$.
Put $a=w_b$ and $r=p_b$. The log-sum inequality on all other categories gives
\[
 \sum_{j\ne b}\bar w_j\log\frac{\bar w_j}{p_j}
 \ge(1-a)\log\frac{1-a}{1-r},
\]
so $\kl(a,r)\le D$. This is data processing under the binary partition
$\{b\},\{b\}^{\rm c}$. On $r\in(0,a)$ the derivative of $\kl(a,r)$ is
$(r-a)/(r(1-r))<0$; the divergence tends to infinity at zero and equals zero
at $r=a$. This proves the inverse bound. Equality is attained by taking
$p_b=r$ and $p_j=(1-r)w_j/(1-a)$ on the other banked categories, with zero
mass off the bank. If strict full support is required, this is the limiting
construction; the infimum is approached as off-bank mass vanishes. The case
$a=1$ is immediate.

Writing $s=p(B)$ and $\widetilde p_b=p_b/s$ gives
$\KL(\bar w\|p)=-\log s+\KL(w\|\widetilde p)\ge-\log s$.
For the coordinate Pinsker bound, fix $r$ and differentiate the binary
KL twice in its first argument: $\partial_a^2\kl(a,r)=1/[a(1-a)]\ge4$.
Its value and first derivative vanish at $a=r$, so
$\kl(a,r)\ge2(a-r)^2$. Finally,
$\kl(a,r)=-h_2(a)-a\log r-(1-a)\log(1-r)
\ge-h_2(a)-a\log r$, giving the elementary exponential bound.
\end{proof}

For the paper's energy inequality, substitute
$C=R_w(p(T))+[\Psi(1)-\Psi(P(T))]/\rho_w$.
This improves the existing termwise bound $p_b\ge e^{-C/w_b}$ because it
uses probability normalization and subtracts target entropy. It can also be
used at an individual checkpoint if a genuine upper bound on its complete
bank cross entropy is available; a checkpoint calculation does not prove
that the same bound holds at later times.

For a uniform bank of 16, $C=\log16+0.01$ certifies
$p_b\ge0.0339712$, whereas the termwise bound is approximately
$4.62\times10^{-20}$. These are illustrative mathematical inputs, not
measured training losses. At excess cross entropy $D=160$, even the sharp
log lower bound is about $-2563.74$, so normalization cannot rescue an
intrinsically enormous energy budget.

\subsection{What a scalar entropy value can certify}
\begin{proposition}[Sharp entropy threshold for all-mode retention]
Let $q$ be a probability vector on exactly $m\ge2$ possible modes. A scalar
bound $H(q)\ge h_0$, with $h_0\le\log m$, gives a strictly positive uniform
coordinate floor if and only if $h_0>\log(m-1)$. In that case
\[
 q_b\ge r_m(h_0)\quad\hbox{for every }b,
\]
where $r_m(h_0)\in(0,1/m]$ uniquely solves
\[
 h_2(r)+(1-r)\log(m-1)=h_0.
\]
The certificate is sharp.
\end{proposition}

\begin{proof}
Fix $r=q_b$. Splitting off that mode gives
\[
 H(q)=h_2(r)+(1-r)H(q_{-b}/(1-r))
 \le h_2(r)+(1-r)\log(m-1)=:f_m(r).
\]
Equality is attained when the other modes are uniform. On $0<r<1/m$,
$f'_m(r)=\log[(1-r)/(r(m-1))]>0$, with
$f_m(0)=\log(m-1)$ and $f_m(1/m)=\log m$.
If $H(q)\ge h_0>\log(m-1)$, a coordinate below the stated root is
impossible. Conversely, if $h_0\le\log(m-1)$, a distribution with one
zero coordinate and the rest uniform satisfies the bound. Even restricting
to positive coordinates permits arbitrarily small values by approximation.
\end{proof}

For uniform $u$, $H(q)=\log m-\KL(q\|u)$, so the same certificate follows
from $\kl(q_b,1/m)\le\KL(q\|u)$. It is positive precisely when the
forward-KL budget is below $\log[m/(m-1)]$. A missing coordinate can have
$H(q)=\log(m-1)$ and $\KL(q\|u)=\log[m/(m-1)]$, while
$\KL(u\|q)=\infty$. For $m=16$, the entropy threshold is $\log15$
(about $2.70805$ nats), only $0.06454$ below $\log16$.
This distinguishes level sets. It does not say that an exact-entropy gradient
flow must reach a zero coordinate; exact entropy can prevent extinction
under suitable dynamics. An empirical entropy estimate with missing modes
is not an exact value of $H(q)$, and conditional entropy alone does not
bound total correct or bank mass.

\subsection{When response likelihood yields a key certificate}
Fix one prompt and finitely many distinct complete responses $e_j$, verified
to produce keys $b_j$. Put a normalized distribution $w_j$ on these responses,
and let $W_b=\sum_{j:b_j=b}w_j$. Under a fixed decoding law, the response
events are disjoint. Extend the exemplar target by zero to every other
complete response in the policy alphabet. Execution is then a deterministic
map on this full alphabet; unlisted responses remain distinguishable until
they are mapped to their keys. Therefore
\[
 -\sum_jw_j\log\pi(e_j\mid x)\le C
 \quad\Longrightarrow\quad
 \kl(W_b,p(b\mid x))\le C-H(w).
\]
Indeed, the response-level KL is $C-H(w)$ or smaller, and deterministic
execution followed by the binary key indicator is data processing.
Thus $p(b\mid x)\ge\ell_-(W_b,C-H(w))$. Individually,
$p(b_j\mid x)\ge\pi(e_j\mid x)$ because the complete-response event is
contained in its verified-key event.

The likelihood must include EOS or another specified termination event,
or a fixed deterministic horizon, and must use the same prompt, temperature,
vocabulary mask, token-truncation rule, and verifier as the probability being
bounded. A scored prefix is not sufficient: its continuations can produce
other keys or invalid outputs. A mean-token log score is not a probability;
for a fixed complete length $L_j$, the sequence probability is
$\exp(L_j\bar\ell_j)$, not $\exp(\bar\ell_j)$.
A sum of length-normalized scores first requires normalization:
if $a_j=\alpha_j/L_j$, $A=\sum_j a_j$, and $w_j=a_j/A$, then
$\sum_j\alpha_j[-\log\pi(e_j)]/L_j=A R_w$.
Responses from different prompts do not share one conditional probability
alphabet; apply this construction prompt by prompt or specify a joint prompt
mixture explicitly. Approximate numerical log scores additionally require
controlled numerical error before being called rigorous certificates.

\subsection{Finite samples distinguish detection from extinction}
At a frozen checkpoint, for a mode specified before sampling, $N$
independent draws give $X\sim\operatorname{Binomial}(N,p_b)$.
Exact binomial-tail inversion supplies one-sided confidence bounds
\cite{clopper1934}. If $X=0$, the one-sided $(1-\alpha)$ upper bound is
\[
 U_0=1-\alpha^{1/N},
\]
with lower bound zero. This follows directly from
$\Prb_{p_b}(X=0)=(1-p_b)^N$: for any $p_b>U_0$, the all-zero event has
probability below $\alpha$. A full upper-confidence procedure uses the
corresponding binomial-tail inversion for nonzero counts as well.

Zero sightings in 32 draws allow $p_b$ as large as $0.08937$ at the
one-sided 95\% level. Simultaneous coverage over 16 prespecified modes by
Bonferroni uses $\alpha/16$ and increases that upper bound to $0.16495$.
At least 299 zero-count draws are needed to obtain a single-mode 95\%
upper bound below $.01$; 574 suffice for simultaneous 16-mode coverage.
These are probability-threshold statements, not proofs of exact extinction.
For example, two sightings in 32 draws give a one-sided 95\% lower bound
of about $.01122$ for one prespecified mode.

With repeated inspection of a fixed policy, an anytime-valid confidence
sequence can avoid invalid optional stopping \cite{howard2021}. One simple
Bernoulli construction uses a fixed beta mixture, with shape parameters
$a,b>0$ and cumulative successes $S_t$:
\[
 M_t(p)=\frac{\mathrm B(S_t+a,t-S_t+b)}
 {\mathrm B(a,b)p^{S_t}(1-p)^{t-S_t}},\qquad
 C_t=\{p:M_t(p)<1/\alpha\}.
\]
For each true $p\in(0,1)$, this is an integral of likelihood-ratio martingales,
so it is a nonnegative mean-one martingale. The maximal inequality for such
martingales gives
$\Prb_p(\exists t:p\notin C_t)\le\alpha$.
Mode-wise error allocations can make this simultaneous over a fixed roster.
The argument requires a fixed Bernoulli success probability; pooling across
changing training checkpoints does not certify the current checkpoint's
probability. Nor do repeated deterministic teacher-forced scores supply
Bernoulli trials.

For an unknown catalogue, missing mass is the total probability of as-yet
unobserved categories under a fixed iid sampling law, not their number and
not evidence that they are extinct. Specialized concentration results address
that total mass \cite{mcallester2003,berend2013}. They cannot establish the
survival of each named banked key, and stationary species formulas cannot be
applied directly to a history generated by changing policies. Named-key
binomial counts are the direct measurement for per-mode visibility.



\section{An exact optimizer comparison from information geometry}
\label{sec:natural-gradient}
Natural gradient on the categorical simplex uses the Fisher metric, also
called the Shahshahani metric in evolutionary dynamics
\cite[Sections~2.3 and 3.1]{harper2009}. This gives a useful comparison with
the paper's Euclidean-logit collapse result.

\begin{proposition}[Natural-gradient selection and replay]
\label{prop:natural-replay}
Start from an interior categorical distribution $p(0)$ with
$0<P(0)<1$. Let $\Psi\in C^2([0,1])$ satisfy
$c(P)=\Psi'(P)>0$, let $w$ be a fixed probability vector supported on a
nonempty correct-mode bank, and extend it by zero to $\bar w$ outside that
bank. Exact Fisher natural-gradient ascent of
\[
 \mathcal J(p)=\Psi(P)+\rho\sum_b w_b\log p_b
\]
with $\rho\ge0$ induces
\begin{align}
 \dot p_i&=c(P)p_i(\mathbf1\{i\in\mathcal C\}-P)
             +\rho(\bar w_i-p_i),\label{eq:natural-p}\\
 \dot P&=(1-P)(c(P)P+\rho),&
 \dot q_c&=\frac{\rho}{P}(w_c-q_c),\label{eq:natural-Pq}
\end{align}
where $w_c=0$ for unbanked correct modes. With $\rho=0$, correctness tends
to one and $q(t)=q(0)$: there is no within-correct-mode collapse.
With $\rho>0$, define $A(t)=\int_0^t\rho/P(s)\,ds$. Then
\[
 q(t)=w+(q(0)-w)e^{-A(t)},\qquad A(t)\ge\rho t,
\]
and
\[
 1-P(t)\le(1-P(0))e^{-\rho t},\qquad
 p_b(t)\ge P(0)\min\{q_b(0),w_b\}>0
\]
for every banked mode and all $t\ge0$.
\end{proposition}
\begin{proof}
For a smooth objective on the interior simplex, the Fisher natural-gradient
vector is $\dot p_i=p_i(\partial_i\mathcal J-
\sum_jp_j\partial_j\mathcal J)$. Here
$\partial_i\mathcal J=c(P)\mathbf1\{i\in\mathcal C\}+
\rho\bar w_i/p_i$, whose $p$-weighted mean is $c(P)P+\rho$.
This proves~\eqref{eq:natural-p}. Summing over correct categories and
differentiating $q_c=p_c/P$ proves~\eqref{eq:natural-Pq}.
Incorrect coordinates obey $\dot p_i=-(c(P)P+\rho)p_i$, so remain positive
at finite times. Correct coordinates stay positive by the explicit
$q$ formula (or are constant in $q$ when $\rho=0$). Thus the interior
calculation applies for all finite times.

For $\rho=0$, $P$ increases and cannot have a limit below one because
$c>0$ on the compact interval $[P(0),1]$. For $\rho>0$, solve the linear
$q$ equation and use $P\le1$ to obtain $A(t)\ge\rho t$.
The differential inequality $\dot P\ge\rho(1-P)$ gives the correctness
bound. Each $q_b(t)$ is between $q_b(0)$ and $w_b$, and $P(t)\ge P(0)$,
which gives the probability floor.
\end{proof}

The term $\rho(\bar w-p)$ has the form of replenishment from a fixed
population distribution. The result is an elementary specialization of the
Fisher-gradient formula, not a claim of a new general evolutionary-dynamics
theorem. It shows that the update geometry changes both collapse and the
strength of retention guarantees. For an incomplete bank it still sends
unbanked correct mass to zero. It supplies no guarantee for approximate
natural gradient, neural Fisher matrices, PPO, or AdamW, and the time variable
is an idealized optimization clock rather than elapsed compute time.


\section{What is ready to use in the paper}
The sharp KL certificate and the exact-entropy distinction directly strengthen
the existing retention discussion. The optimizer theorem extends the fixed
joint-potential argument to controlled stochastic updates. The admission
propositions identify sufficient conditions for a discovery-to-retention
claim, including the effects of finite capacity and bank changes.
The natural-gradient calculation is an optional mechanism comparison that
makes the role of optimizer geometry particularly explicit.

Applying these statements to the released neural runs still requires their
premises to be measured or established. In particular, exact complete-response
likelihood and finite categorical replay cross entropy must not be confused
with mean-token scores. A theorem about probability retention cannot substitute
for observations that individual stored modes survived the actual training
trajectory. The registered E121 study remains prospective at this audit:
no per-mode telemetry was available in its five registered run directories.
The accompanying survival audit identifies analysis work needed before those
records can support the registered claims.

\begin{thebibliography}{99}
\bibitem{pereyra2017} Gabriel Pereyra, George Tucker, Jan Chorowski,
\L ukasz Kaiser, and Geoffrey Hinton (2017).
\emph{Regularizing Neural Networks by Penalizing Confident Output Distributions}.
\url{https://arxiv.org/abs/1701.06548}, Section~3.2.
\bibitem{agarwal2021} Alekh Agarwal, Sham M. Kakade, Jason D. Lee,
and Gaurav Mahajan (2021).
\emph{On the Theory of Policy Gradient Methods: Optimality, Approximation,
and Distribution Shift}. JMLR 22(98):1--76.
\url{https://jmlr.org/papers/v22/19-736.html}, Section~5.2.
\bibitem{geist2019} Matthieu Geist, Bruno Scherrer, and Olivier Pietquin (2019).
\emph{A Theory of Regularized Markov Decision Processes}.
ICML, PMLR 97:2160--2169.
\url{https://proceedings.mlr.press/v97/geist19a.html}.
\bibitem{mei2020} Jincheng Mei, Chenjun Xiao, Csaba Szepesv\'ari,
and Dale Schuurmans (2020).
\emph{On the Global Convergence Rates of Softmax Policy Gradient Methods}.
ICML, PMLR 119:6820--6829.
\url{https://proceedings.mlr.press/v119/mei20b.html}.
\bibitem{harper2009} Marc Harper (2009).
\emph{Information Geometry and Evolutionary Game Theory}.
\url{https://arxiv.org/abs/0911.1383}, Sections~2.3 and~3.1.
\bibitem{durrett2019} Rick Durrett (2019).
\emph{Probability: Theory and Examples}, fifth edition.
Cambridge University Press. Author-hosted manuscript:
\url{https://sites.math.duke.edu/~rtd/PTE/PTE5_011119.pdf}, Theorem~4.3.4.
\bibitem{branicky1998} Michael S. Branicky (1998).
\emph{Multiple Lyapunov Functions and Other Analysis Tools for Switched and Hybrid Systems}.
IEEE Transactions on Automatic Control 43(4):475--482.
\url{https://people.ece.ubc.ca/moishi/eece571m/articles/Branicky98.pdf}.
\bibitem{ecoffet2021} Adrien Ecoffet, Joost Huizinga, Joel Lehman,
Kenneth O. Stanley, and Jeff Clune (2021).
\emph{First return, then explore}. Nature 590:580--586.
\url{https://doi.org/10.1038/s41586-020-03157-9}.

\bibitem{vanerven2014} Tim van Erven and Peter Harremo\"es (2014).
\emph{R\'enyi Divergence and Kullback--Leibler Divergence}.
IEEE Transactions on Information Theory, DOI
\href{https://doi.org/10.1109/TIT.2014.2320500}{10.1109/TIT.2014.2320500}.
Verified author version: \href{https://arxiv.org/abs/1206.2459v2}{arXiv:1206.2459v2};
Theorems 9 and 31, and Equation (2).
\bibitem{garivier2011} Aur\'elien Garivier and Olivier Capp\'e (2011).
\emph{The KL-UCB Algorithm for Bounded Stochastic Bandits and Beyond}.
Proceedings of COLT, PMLR 19:359--376.
\url{https://proceedings.mlr.press/v19/garivier11a.html}.
\bibitem{clopper1934} C. J. Clopper and E. S. Pearson (1934).
\emph{The Use of Confidence or Fiducial Limits Illustrated in the Case of the Binomial}.
Biometrika 26(4):404--413.
\href{https://doi.org/10.1093/biomet/26.4.404}{doi:10.1093/biomet/26.4.404}.
\bibitem{howard2021} Steven R. Howard, Aaditya Ramdas, Jon McAuliffe, and Jasjeet Sekhon (2021).
\emph{Time-uniform, nonparametric, nonasymptotic confidence sequences}.
Annals of Statistics 49(2):1055--1080.
\href{https://doi.org/10.1214/20-AOS1991}{doi:10.1214/20-AOS1991};
\href{https://arxiv.org/abs/1810.08240}{arXiv:1810.08240}, Proposition 7.
\bibitem{mcallester2003} David McAllester and Luis Ortiz (2003).
\emph{Concentration Inequalities for the Missing Mass and for Histogram Rule Error}.
Journal of Machine Learning Research 4:895--911.
\url{https://jmlr.org/papers/volume4/mcallester03a/mcallester03a.pdf}.
\bibitem{berend2013} Daniel Berend and Aryeh Kontorovich (2013).
\emph{On the concentration of the missing mass}.
Electronic Communications in Probability 18, paper 3:1--7.
\href{https://doi.org/10.1214/ECP.v18-2359}{doi:10.1214/ECP.v18-2359};
\href{https://arxiv.org/abs/1210.3248}{arXiv:1210.3248}.

\end{thebibliography}
\end{document}
